\myparagraph{Types}
Let $\vocab$ be a vocabulary with unary  $U_1, \ldots, U_n$ and binary  $R_1, \ldots, R_m$.
\begin{definition}[1- and 2-type]
    A \emph{1-type} (w.r.t. $\vocab$) is a maximal consistent subset of
    \begin{equation*}
        \bigcup_{i \in \intsinterval{n}}\set{U_i(x), \neg U_i(x)}
        \ \cup\ 
        \bigcup_{i \in \intsinterval{m}}\set{R_i(x, x), \neg R_i(x, x)}.
    \end{equation*}
    A \emph{2-type} (w.r.t. $\vocab$) is a maximal consistent subset of
    \begin{equation*}
        \bigcup_{i \in \intsinterval{m}}
        \set{R_i(x, y), \neg R_i(x, y), R_i(y, x), \neg R_i(y, x)}.
    \end{equation*}
\end{definition}
We denote the set of $1$-types  by $\onetypes$ and the set of $2$-types $\twotypes$.

Given an element $v$ in a $\vocab$-structure $\cG$, we let $\onetypeof(v)$ denote its one-type in $\cG$
and similarly for a pair of elements $v, u$ in $\cG$ we let $\twotypeof(v,u)$ denote its two-type in $\cG$;
in both cases we omit the dependence on $\cG$ for brevity.

The following definitions are variations of notions
in Pratt-Hartmann's~\cite{ctwocomplexity,gctwocomplexity,ianrevisited}.

\begin{definition}[Silent 2-type]
    The \emph{silent type}, denoted by $\etanull$, is the unique 2-type that consists of only negated terms.
    If a type is not silent we say it is \emph{audible}.
\end{definition}

\begin{definition}[Dual of 2-type]
    For a 2-type $\eta$, the \emph{dual} of $\eta$, denoted by $\dual{\eta}$,
    it the unique 2-type that swaps $x$ and $y$ in $\eta$.
\end{definition}
We denote
the set of all 2-types except $\etanull$ by $\twotypes^+$, and
the set of all 2-types that contain $R_t(x, y)$ by $\twotypes_t$ for $1 \le t \le m$.

Note that $\abs{\Onetypes} = 2^{n+m}$,
$\abs{\twotypes} = 2^{2m}$, $\abs{\twotypes^+} = 2^{2m}-1$,
and $\abs{\twotypes_t} =2^{2m-1}$.

We can abstract a model based on types:
\begin{definition}[Type graph]
    We associate to a $\vocab$-structure in a binary vocabulary a colored graph in an exponentially larger vocabulary:
    each vertex is colored with a 1-type, and
    each edge is colored with a 2-type. This is the \emph{type graph} of the $\vocab$-structure.
\end{definition}

Losing information, we can abstract a finite model based on cardinalities of $1$-types. The notion will be crucial in Section \ref{sec:ctwo}:
\begin{definition}[$1$-type cardinality vector of a $\vocab$-structure] \label{def:cardboundvector}
    Let $\cG$ be a $\vocab$-structure.
    The $1$-type cardinality vector of $\cG$ is a $\abs{\onetypes}$-dimensional vector 
    where the $i^{th}$ component is the number of elements in $\cG$ with 1-type $\pi_i$.
\end{definition}

